% =====================================================================
%  Chapter 0: Preface — Caveat Computor
%  Source: Intro.tex (Overleaf) — adapted and expanded
% =====================================================================

\chapter*{Preface: \textit{Caveat Computor}}\label{chap_Preface}
\addcontentsline{toc}{chapter}{Preface: \textit{Caveat Computor}}

\section*{What this book is about}

The primary objective of this manuscript is to introduce students to the
quantitative methods and techniques used in urban and regional research. The
emphasis is on using methods in the context of planning, urban policy, and
regional science research: matching the method to the problem, developing an
interactive engagement with data, promoting creative and accessible styles of
data presentation, and cultivating a critical literacy of methods and their
limitations.

In its intellectual lineage, this text is unambiguously in the spirit of
Walter \citeapos{Isard1960a} classic \emph{Methods of Regional Analysis}---a
methodological manifesto for regional science as a scientific project, if ever
there was one.\footnote{Now in its seventh edition, \emph{Methods of
Interregional and Regional Analysis} \citep{Isard1998} remains the
methodological benchmark for regional science, covering regional and
interregional input-output analysis, econometrics, programming and industrial
and urban complex analysis, gravity and spatial interaction models, and applied
general interregional equilibrium models. In contrast to the largely
theoretical treatment in that volume, this manual focuses on the empirical
aspects of application, with special emphasis on the connections between theory
and evidence.} Of all the subdisciplines subsumed under the heading of
regional science, urban and regional planning is perhaps the archetypical
interdisciplinary venture---yet it still relies on solid disciplinary
skills. Planning educators regularly define an accomplished planner as
simultaneously ``the best geographer, the best economist, the best
environmental scientist, the best political scientist, and the best
architect'' \citep{Fainstein2011}.

\section*{The computational abundance paradox}

As computational costs have fallen precipitously and data availability has
exploded, one might expect the quality of empirical analysis to have risen in
proportion. The opposite tendency is equally plausible---and, this author
believes, equally common. Software and analytical tools are becoming ever more
user-friendly, and hitherto prohibitively complex analyses now require only a
few clicks deployed via an attractive graphical user interface. Yet perhaps
paradoxically, the more we rely on technology and the more accessible data
tools become, the less time is spent on careful and critical analysis. This
leads to a latent tendency to produce more empirical output rather than better
analysis. Just because we can (cheaply and conveniently) compute something does
not mean that we should. In other words: \emph{caveat computor}.

This paradox has only deepened with the emergence of large language model
(LLM) coding assistants capable of generating spatial regression code from a
single natural-language prompt. When generating a spatial lag model requires
less effort than reading this sentence, the limiting constraint on analysis
quality shifts entirely to whether the analyst understands what the model is
doing---whether the weight matrix choice is defensible, whether the
identification strategy is credible, whether the coefficient can be given a
causal interpretation. The theoretical rigour of Parts~I and~II of this
book is \emph{justified by}, not in tension with, the ease of modern code
generation. The most dangerous analyst is one who can produce polished output
without understanding its assumptions.

\section*{A new paradigm: regional policy informatics}

This book proposes a synthesis that takes the credibility revolution in
econometrics seriously, treats spatial structure as an inferential challenge
rather than a nuisance parameter, and integrates the quantitative traditions
of urban economics, regional science, economic geography, and real estate
around the unifying framework of spatial equilibrium.

\noindent\textbf{Central claim.} Spatial analysis in the social sciences has
systematically produced \emph{geographically-correlated} answers rather than
\emph{causally-identified} ones. Geographic confounding---the systematic
confusion of spatial correlation with causal identification---is to regional
analytics what omitted variable bias is to cross-sectional econometrics. The
response is a credibility-first, spatially-literate workflow built entirely on
open-source tools and public data infrastructure.

We define \textbf{regional policy informatics} as the systematic application
of computational, statistical, and causal inference methods to the analysis of
spatially-structured policy problems, using reproducible workflows and open
data infrastructure. What distinguishes this from ``spatial data science''
(the term used by the Python-oriented GIS community) is the explicit
commitment to policy-relevant causal identification. Spatial data science
optimises for pattern recognition and visualisation---it is largely descriptive
and predictive. Regional policy informatics is evaluative: it asks whether
interventions work, for whom, and why.

\section*{Four questions that frame good quantitative research}

The analytical philosophy underlying this text is organised around four
questions that \citet{Angrist2009a} argue should frame any credible empirical
project:

\begin{enumerate}
  \item What is the \emph{causal relationship of interest}? This question
    gives rise to the importance of the counterfactual.
  \item What \emph{ideal experiment} would capture this causal effect?
  \item What is the \emph{identification strategy}? How is the observational
    data used to approximate the ideal experiment?
  \item What is the \emph{mode of inference}? (microdata, clustered groups,
    spatial units, etc.)
\end{enumerate}

These four questions are not only the organising principle of Part~II
(The Linear Model and Causal Inference); they reappear throughout the book
as a touchstone for evaluating any spatial or regional analysis.

\section*{The open-source commitment}

This text is built around three principles:

\begin{enumerate}
  \item \textbf{Open source software} achieves interoperability and
    reproducibility without licensing barriers. The full analytical toolkit---
    from spatial regression to input-output modelling---is implemented in R
    and Python using freely available packages.

  \item \textbf{Public data infrastructure} supports a full regional
    analytics workflow without proprietary data. The Census, ACS, FRED,
    LEHD, BEA Regional Accounts, HUD, and OpenStreetMap collectively provide
    richer inputs than most proprietary alternatives.

  \item \textbf{Replicable research} is a scientific norm, not a technical
    nicety. The \emph{American Economic Review} requires that all data and
    code sufficient for replication be deposited prior to publication
    \citep{AER2024}. Urban planning and geography journals have been slower to
    adopt such standards; this text argues that they should not be.
\end{enumerate}

\section*{Prerequisites and how to use this book}

The material is aimed at first-year graduate students (or advanced
undergraduates) who have taken at minimum an introductory statistics course
covering the basics of probability, confidence intervals, hypothesis testing,
and bivariate regression. Exposure to linear algebra is helpful but not
required---Part~I provides everything needed from scratch.

Chapter~1 reviews probability and statistics. Chapter~2 introduces the matrix
algebra toolkit that underlies all subsequent chapters. Chapter~3 surveys data
sources and measurement. These three chapters constitute Part~I and can be
read selectively by students with stronger preparation. The heart of the book
is Parts~II and~III: Chapters~4--6 develop the linear model and causal
inference toolkit; Chapters~7--9 apply that toolkit to the spatial economy.
Part~IV covers time series, panel data, and the frontiers of computational
regional analysis.

\renewcommand\bibname{Further reading}
\begin{subappendices}
\bibliographystyle{econometrica}
\bibliography{Literature/OverLord}
\IfFileExists{Literature/Preface.bbl}{\input{Literature/Preface.bbl}}{\typeout{Need to run bibtex on Preface}}
\end{subappendices}
